The development of robust clinical decision support systems is frequently impeded by the scarcity of high-fidelity biomedical data \rev{that can be shared with limited disclosure risk}. While Generative Large Language Models (LLMs) offer a promising avenue for synthetic data generation, they often struggle to capture the complex, non-linear dependencies and severe class imbalances inherent in Electronic Health Records (EHR), leading to statistically plausible but clinically invalid records. To bridge this gap, we introduce DISCO-TAB (DIScriminator-guided COntrol for TABular synthesis), a novel framework that orchestrates a fine-tuned LLM with a multi-objective discriminator system optimized via Reinforcement Learning. Unlike prior methods relying on scalar feedback, DISCO-TAB evaluates synthesis at four granularities, token, sentence, feature, and row, while integrating Automated Constraint Discovery and Inverse-Frequency Reward Shaping to autonomously preserve latent medical logic and resolve minority-class collapse. We \rev{evaluate} our framework across \rev{nine} benchmarks, including high-dimensional, small-sample medical datasets (e.g., Heart Failure, Parkinson's). \rev{Our results demonstrate that hierarchical feedback yields strong downstream utility, improving clinical classifier F1 by up to $+0.382$ in absolute terms (a 38.2 percentage-point gain) over the strongest GAN, diffusion, and LLM baselines, while maintaining low distributional divergence (per-dataset Jensen--Shannon divergence, JSD, reported in Table~\ref{tab:unified_fidelity_baseline}) and favorable empirical, proximity-based disclosure-risk indicators. We emphasize that DISCO-TAB does not implement a formal differential-privacy mechanism; its privacy behavior is assessed empirically rather than guaranteed. This work contributes a practical framework for generating utility-preserving synthetic tabular data to support sensitive healthcare applications.}